\section{Ação de melhoria}
\label{sec:melhoria}

O roteiro pede uma ação de melhoria justificada pelos resultados \cite{ramos2026roteiro}. Como o SIGAA é um produto de código fechado mantido pela STI/UnB, a equipe escolheu um \textbf{protótipo navegável de alta fidelidade}, \textit{mobile-first} (360~px), das tarefas T1 a T6. A escolha segue o julgamento: Capacidade de Interação é a característica de pior resultado, T6 e T4 são as tarefas de pior índice, e a falta de adaptação ao celular (M14 = 0\%) atinge os cinco participantes, que usam o SIGAA no celular.

O protótipo implementa dez requisitos, cada um ligado a uma evidência da Fase 4: declaração de \textit{viewport} e \textit{layout} em coluna única (M14); carga horária pendente e percentual integralizado no topo da tela inicial (T4); horário por extenso ao lado do código (M11); barra inferior com quatro abas e botões de pelo menos 46~px, sem menus por passagem do mouse (M10 e critério (e) de M14); calendário acadêmico vigente dentro do SIGAA, em um toque, no lugar do link para a página da Pós-Graduação (T6); uma central única de materiais na Turma Virtual (T5); idioma, textos alternativos, rótulos e contraste AA (M9); grade semanal das turmas (M3); carregamento sob demanda (dados de desempenho); e preservação, sem ação, das funções fora do escopo, para manter o mapa do portal real.

\textbf{Verificação.} O protótipo foi medido com os mesmos critérios e \textit{scripts} aplicados ao SIGAA: o \textit{script} de M14 levou-o ao estado final de cada tarefa, por toques, em 360 e 768~px, e o motor do WAVE (versão 3.3.4.0, a mesma de M9) avaliou cada estado. A primeira medição deu M14 = 0 de 12, com textos de 9,5 a 11,5~px e botões de 22~px; o protótipo foi ajustado (fonte mínima de 12~px e alvos de pelo menos 24~px) e medido de novo. O resultado foi \textbf{M14 = 12 de 12 (Excelente)} e \textbf{M9 = 0 erro por página (Excelente)}, contra 0 de 12 e 13,17 no SIGAA, e T6 passou a ser feita em um toque. O efeito sobre o tempo, a conclusão e a satisfação não foi medido, porque não houve novo teste com usuários.
